\chapter{2021年全国甲卷（文）}
\sloppy
\hfuzz=120pt

\section{单选题}

\begin{problem}
设集合 \(M=\{1,3,5,7,9\}\)，\(N=\{x\mid 2x>7\}\)，则 \(M\cap N=\)（\quad）

\choices
    {\(\{7,9\}\)}
    {\(\{5,7,9\}\)}
    {\(\{3,5,7,9\}\)}
    {\(\{1,3,5,7,9\}\)}

\end{problem}

\begin{answer}
B
\end{answer}

\begin{solution}
由 \(2x>7\) 得 \(x>\frac{7}{2}\)，所以 \(N=\left(\frac{7}{2},+\infty\right)\).
故 \(M\cap N=\{5,7,9\}\).
故选 B.
\end{solution}

\begin{problem}
为了解某地农村经济情况，对该地农户家庭年收入进行抽样调查，将农户家庭年收入的调查数据整理得到如下频率分布直方图：

\texfigure[max width=4.6cm,max totalheight=4.8cm]{img_repaint/2021/national_paper_a_liberal/q2_fig1.tex}

根据此频率分布直方图，下面结论中不正确的是（\quad）

\choices
    {该地农户家庭年收入低于 \(4.5\) 万元的农户比率估计为 \(6\%\)}
    {该地农户家庭年收入不低于 \(10.5\) 万元的农户比率估计为 \(10\%\)}
    {估计该地农户家庭年收入的平均值不超过 \(6.5\) 万元}
    {估计该地有一半以上的农户，其家庭年收入介于 \(4.5\) 万元至 \(8.5\) 万元之间}

\end{problem}

\begin{answer}
C
\end{answer}

\begin{solution}
因为频率直方图中的组距为 \(1\)，所以各组小矩形的高等于频率，样本频率可以作为总体相应比率的估计值.

该地农户家庭年收入低于 \(4.5\) 万元的农户比率估计值为 \(0.02+0.04=0.06=6\%\).
故 \(A\) 正确.

该地农户家庭年收入不低于 \(10.5\) 万元的农户比率估计值为 \(0.04+0.02\times 3=0.10=10\%\).
故 \(B\) 正确.

该地农户家庭年收入介于 \(4.5\) 万元至 \(8.5\) 万元之间的比率估计值为 \(0.10+0.14+0.20\times 2=0.64=64\%>50\%\).
故 \(D\) 正确.

该地农户家庭年收入平均值的估计值为
\examdisplaychain{
&3\times 0.02+4\times 0.04+5\times 0.10+6\times 0.14+7\times 0.20+8\times 0.20\\
&\quad+9\times 0.10+10\times 0.10+11\times 0.04+12\times 0.02+13\times 0.02+14\times 0.02\\
=&7.68
}
（万元），超过 \(6.5\) 万元，故 \(C\) 错误.
综上，给出结论中不正确的是 \(C\).
\end{solution}

\begin{problem}
已知 \((1-\i)^2z=3+2\i\)，则 \(z=\)（\quad）

\choices
    {\(-1-\frac{3}{2}\i\)}
    {\(-1+\frac{3}{2}\i\)}
    {\(-\frac{3}{2}+\i\)}
    {\(-\frac{3}{2}-\i\)}

\end{problem}

\begin{answer}
B
\end{answer}

\begin{solution}
由 \((1-\i)^2=-2\i\)，得 \(z=\frac{3+2\i}{-2\i}=\frac{(3+2\i)\i}{-2\i^2}=\frac{-2+3\i}{2}=-1+\frac{3}{2}\i\).
故选 B.
\end{solution}

\begin{problem}
下列函数中是增函数的是（\quad）

\choices
    {\(f(x)=-x\)}
    {\(f(x)=\left(\frac{2}{3}\right)^x\)}
    {\(f(x)=x^2\)}
    {\(f(x)=\sqrt[3]{x}\)}

\end{problem}

\begin{answer}
D
\end{answer}

\begin{solution}
对于 \(A\)，函数 \(f(x)=-x\) 是 \( \R \) 上的减函数，不合题意.

对于 \(B\)，函数 \(f(x)=\left(\frac{2}{3}\right)^x\) 是 \( \R \) 上的减函数，不合题意.

对于 \(C\)，函数 \(f(x)=x^2\) 在 \((-\infty,0)\) 上是减函数，不合题意.

对于 \(D\)，函数 \(f(x)=\sqrt[3]{x}\) 是 \( \R \) 上的增函数，符合题意.

故选 D.
\end{solution}

\begin{problem}
点 \((3,0)\) 到双曲线 \(\frac{x^2}{16}-\frac{y^2}{9}=1\) 的一条渐近线的距离为（\quad）

\choices
    {\(\frac{9}{5}\)}
    {\(\frac{8}{5}\)}
    {\(\frac{6}{5}\)}
    {\(\frac{4}{5}\)}

\end{problem}

\begin{answer}
A
\end{answer}

\begin{solution}
由题意可知，双曲线的渐近线方程为 \(\frac{x^2}{16}-\frac{y^2}{9}=0\)，即 \(3x\pm 4y=0\).
结合对称性，不妨考虑点 \((3,0)\) 到直线 \(3x+4y=0\) 的距离，则
\(d=\frac{|3\times 3+4\times 0|}{\sqrt{3^2+4^2}}=\frac{9}{5}\).
故选 A.
\end{solution}

\begin{problem}
青少年视力是社会普遍关注的问题，视力情况可借助视力表测量. 通常用五分记录法和小数记录法记录视力数据，五分记录法的数据 \(L\) 和小数记录法的数据 \(V\) 满足 \(L=5+\lg V\). 已知某同学视力的五分记录法的数据为 \(4.9\)，则其视力的小数记录法的数据为（\quad）（\(\sqrt[10]{10}\approx 1.259\)）

\choices
    {\(1.5\)}
    {\(1.2\)}
    {\(0.8\)}
    {\(0.6\)}

\end{problem}

\begin{answer}
C
\end{answer}

\begin{solution}
由 \(L=5+\lg V\)，当 \(L=4.9\) 时，\(\lg V=-0.1\).
所以 \(V=10^{-0.1}=10^{-\frac{1}{10}}=\frac{1}{\sqrt[10]{10}}\approx \frac{1}{1.259}\approx 0.8\).
故选 C.
\end{solution}

\begin{problem}
在一个正方体中，过顶点 \(A\) 的三条棱的中点分别为 \(E\)，\(F\)，\(G\). 该正方体截去三棱锥 \(A-EFG\) 后，所得多面体的三视图中，正视图如图所示，则相应的侧视图是（\quad）

\bitmapfigure[max width=7cm,max totalheight=4.8cm]{img/2021/national_paper_A_liberal/q7_fig1.png}

\choices
    {\choicebitmap[max width=3.0cm,max totalheight=3.2cm]{img/2021/national_paper_A_liberal/q7_optA.png}}
    {\choicebitmap[max width=3.0cm,max totalheight=3.2cm]{img/2021/national_paper_A_liberal/q7_optB.png}}
    {\choicebitmap[max width=3.0cm,max totalheight=3.2cm]{img/2021/national_paper_A_liberal/q7_optC.png}}
    {\choicebitmap[max width=3.0cm,max totalheight=3.2cm]{img/2021/national_paper_A_liberal/q7_optD.png}}

\end{problem}

\begin{answer}
D
\end{answer}

\begin{solution}
由题意及正视图可得几何体的直观图，恢复后其侧视图对应图 \(D\). 故选 D.
\end{solution}

\begin{problem}
在 \(\triangle ABC\) 中，已知 \(B=120^\circ\)，\(AC=\sqrt{19}\)，\(AB=2\)，则 \(BC=\)（\quad）

\choices
    {\(1\)}
    {\(\sqrt{2}\)}
    {\(\sqrt{5}\)}
    {\(3\)}

\end{problem}

\begin{answer}
D
\end{answer}

\begin{solution}
设 \(AB=c\)，\(AC=b\)，\(BC=a\)，则 \(c=2\)，\(b=\sqrt{19}\).
由余弦定理得 \(b^2=a^2+c^2-2ac\cos B\)，即 \(19=a^2+4-2\times a\times 2\times \cos 120^\circ\).
因为 \(\cos 120^\circ=-\frac{1}{2}\)，所以 \(a^2+2a-15=0\).
解得 \(a=3\) 或 \(a=-5\)（舍去），故 \(BC=3\).
故选 D.
\end{solution}

\begin{problem}
记 \(S_n\) 为等比数列 \(\{a_n\}\) 的前 \(n\) 项和. 若 \(S_2=4\)，\(S_4=6\)，则 \(S_6=\)（\quad）

\choices
    {7}
    {8}
    {9}
    {10}

\end{problem}

\begin{answer}
A
\end{answer}

\begin{solution}
因为 \(S_n\) 为等比数列 \(\{a_n\}\) 的前 \(n\) 项和，所以 \(S_2\)，\(\ S_4-S_2\)，\(\ S_6-S_4\) 成等比数列.
由 \(S_2=4\)，\(S_4=6\)，得 \(S_4-S_2=2\).
所以 \(S_6-S_4=\frac{(S_4-S_2)^2}{S_2}=\frac{2^2}{4}=1\).
于是 \(S_6=S_4+1=7\).
故选 A.
\end{solution}

\begin{problem}
将 \(3\) 个 \(1\) 和 \(2\) 个 \(0\) 随机排成一行，则 \(2\) 个 \(0\) 不相邻的概率为（\quad）

\choices
    {0.3}
    {0.5}
    {0.6}
    {0.8}

\end{problem}

\begin{answer}
C
\end{answer}

\begin{solution}
将 \(3\) 个 \(1\) 和 \(2\) 个 \(0\) 随机排成一行，可以是 \(00111\)，\(\ 01011\)，\(\ 01101\)，\(\ 01110\)，\(\ 10011\)，\(\ 10101\)，\(\ 10110\)，\(\ 11001\)，\(\ 11010\)，\(\ 11100\)，共 \(10\) 种排法.
其中 \(2\) 个 \(0\) 不相邻的排法为 \(01011\)，\(\ 01101\)，\(\ 01110\)，\(\ 10101\)，\(\ 10110\)，\(\ 11010\)，共 \(6\) 种.
所以所求概率为 \(\frac{6}{10}=0.6\).
故选 C.
\end{solution}

\begin{problem}
若 \(\alpha\in\left(0,\frac{\pi}{2}\right)\)，\(\tan 2\alpha=\frac{\cos\alpha}{2-\sin\alpha}\)，则 \(\tan\alpha=\)（\quad）

\choices
    {\(\frac{\sqrt{15}}{15}\)}
    {\(\frac{\sqrt{5}}{5}\)}
    {\(\frac{\sqrt{5}}{3}\)}
    {\(\frac{\sqrt{15}}{3}\)}

\end{problem}

\begin{answer}
A
\end{answer}

\begin{solution}
由二倍角公式，\(\tan 2\alpha=\frac{\sin 2\alpha}{\cos 2\alpha}=\frac{2\sin\alpha\cos\alpha}{1-2\sin^2\alpha}\). 又由题意 \(\tan 2\alpha=\frac{\cos\alpha}{2-\sin\alpha}\). 因为 \(\alpha\in\left(0,\frac{\pi}{2}\right)\)，所以 \(\cos\alpha\neq 0\)，于是 \(\frac{2\sin\alpha}{1-2\sin^2\alpha}=\frac{1}{2-\sin\alpha}\).
解得 \(\sin\alpha=\frac{1}{4}\). 所以 \(\cos\alpha=\sqrt{1-\sin^2\alpha}=\sqrt{1-\frac{1}{16}}=\frac{\sqrt{15}}{4}\).
故 \(\tan\alpha=\frac{\sin\alpha}{\cos\alpha}=\frac{\frac{1}{4}}{\frac{\sqrt{15}}{4}}=\frac{1}{\sqrt{15}}=\frac{\sqrt{15}}{15}\).
故选 A.
\end{solution}

\begin{problem}
设 \(f(x)\) 是定义域为 \( \R \) 的奇函数，且 \(f(1+x)=f(-x)\). 若 \(f\left(-\frac{1}{3}\right)=\frac{1}{3}\)，则 \(f\left(\frac{5}{3}\right)=\)（\quad）

\choices
    {\(-\frac{5}{3}\)}
    {\(-\frac{1}{3}\)}
    {\(\frac{1}{3}\)}
    {\(\frac{5}{3}\)}

\end{problem}

\begin{answer}
C
\end{answer}

\begin{solution}
由题意，\(f\left(\frac{5}{3}\right)=f\left(1+\frac{2}{3}\right)=f\left(-\frac{2}{3}\right)=-f\left(\frac{2}{3}\right)\).
又 \(f\left(\frac{2}{3}\right)=f\left(1-\frac{1}{3}\right)=f\left(\frac{1}{3}\right)=-f\left(-\frac{1}{3}\right)=-\frac{1}{3}\).
所以 \(f\left(\frac{5}{3}\right)=-f\left(\frac{2}{3}\right)=\frac{1}{3}\).
故选 C.
\end{solution}
\section{填空题}

\begin{problem}
若向量 \(\bs a\)，\(\bs b\) 满足 \(|\bs a|=3\)，\(|\bs a-\bs b|=5\)，\(\bs a\cdot\bs b=1\)，则 \(|\bs b|=\fillinblank{}\).

\end{problem}

\begin{answer}
\(3\sqrt{2}\)
\end{answer}

\begin{solution}
由 \(|\bs a-\bs b|=5\)，得 \(|\bs a-\bs b|^2=|\bs a|^2+|\bs b|^2-2\bs a\cdot\bs b\).
代入已知条件，得 \(25=9+|\bs b|^2-2\).
所以 \(|\bs b|^2=18\)，从而 \(|\bs b|=3\sqrt{2}\).
故答案为 \(3\sqrt{2}\).
\end{solution}

\begin{problem}
已知一个圆锥的底面半径为 \(6\)，其体积为 \(30\pi\)，则该圆锥的侧面积为\fillinblank{}.

\end{problem}

\begin{answer}
\(39\pi\)
\end{answer}

\begin{solution}
由圆锥体积公式 \(V=\frac{1}{3}\pi r^2h=30\pi\) 得 \(\frac{1}{3}\pi\times 6^2\times h=30\pi\)，解得 \(h=\frac{5}{2}\).
设母线长为 \(l\)，则 \(l=\sqrt{h^2+r^2}=\sqrt{\left(\frac{5}{2}\right)^2+6^2}=\sqrt{\frac{25}{4}+\frac{144}{4}}=\frac{13}{2}\).
所以圆锥的侧面积为 \(S_{\text{侧}}=\pi rl=\pi\times 6\times \frac{13}{2}=39\pi\).
故答案为 \(39\pi\).
\end{solution}

\begin{problem}
已知函数 \(f(x)=2\cos(\omega x+\varphi)\) 的部分图像如图所示，则 \(f\left(\frac{\pi}{2}\right)=\fillinblank{}\).

\bitmapfigure[max width=6.0cm,max totalheight=4.8cm]{img/2021/national_paper_A_liberal/q15_fig1.png}

\end{problem}

\begin{answer}
\(-\sqrt{3}\)
\end{answer}

\begin{solution}
由题意可得 \(\frac{3}{4}T=\frac{13\pi}{12}-\frac{\pi}{3}=\frac{3\pi}{4}\)，所以 \(T=\pi\)，\(\omega=\frac{2\pi}{T}=2\).
当 \(x=\frac{13\pi}{12}\) 时，图像取得最大值 \(2\)，所以 \(\omega x+\varphi=2k\pi\)，即 \(2\times \frac{13\pi}{12}+\varphi=2k\pi\)，从而 \(\varphi=2k\pi-\frac{13\pi}{6}\). 令 \(k=1\)，得 \(\varphi=-\frac{\pi}{6}\).
所以 \(f(x)=2\cos\left(2x-\frac{\pi}{6}\right)\).
于是 \(f\left(\frac{\pi}{2}\right)=2\cos\left(2\times \frac{\pi}{2}-\frac{\pi}{6}\right)=2\cos\frac{5\pi}{6}=-\sqrt{3}\).
故答案为 \(-\sqrt{3}\).
\end{solution}

\begin{problem}
已知 \(F_1\)，\(F_2\) 为椭圆 \(C:\frac{x^2}{16}+\frac{y^2}{4}=1\) 的两个焦点，\(P\)，\(Q\) 为 \(C\) 上关于坐标原点对称的两点，且 \(|PQ|=|F_1F_2|\)，则四边形 \(PF_1QF_2\) 的面积为\fillinblank{}.

\end{problem}

\begin{answer}
8
\end{answer}

\begin{solution}
因为 \(P\)，\(Q\) 为 \(C\) 上关于坐标原点对称的两点，且 \(|PQ|=|F_1F_2|\)，所以四边形 \(PF_1QF_2\) 为矩形.
设 \(|PF_1|=m\)，\(|PF_2|=n\)，则 \(m+n=2a=8\).
又因为矩形 \(PF_1QF_2\) 的对角线相等，所以 \(\lvert F_1F_2\rvert^2=\lvert PQ\rvert^2=m^2+n^2\).
而 \(\lvert F_1F_2\rvert^2=(2c)^2=4(a^2-b^2)=4(16-4)=48\)，故 \(m^2+n^2=48\).
由 \((m+n)^2=m^2+2mn+n^2\) 得 \(64=48+2mn\)，所以 \(mn=8\).
因此四边形 \(PF_1QF_2\) 的面积为 \(8\).
\end{solution}
\section{解答题}

\begin{problem}
甲，乙两台机床生产同种产品，产品按质量分为一级品和二级品. 为了比较两台机床产品的质量，分别用两台机床各生产了 \(200\) 件产品，产品的质量情况统计如下表：

\begin{center}
\begin{tabular}{|c|c|c|c|}
\hline
 & 一级品 & 二级品 & 合计 \\
\hline
甲机床 & 150 & 50 & 200 \\
\hline
乙机床 & 120 & 80 & 200 \\
\hline
合计 & 270 & 130 & 400 \\
\hline
\end{tabular}
\end{center}

\begin{enumerate}
\item 甲机床，乙机床生产的产品中一级品的频率分别是多少？
\item 能否有 \(99\%\) 的把握认为甲机床的产品质量与乙机床的产品质量有差异？
\end{enumerate}

附：\(K^2=\dfrac{n(ad-bc)^2}{(a+b)(c+d)(a+c)(b+d)}\)

\begin{center}
\begin{tabular}{|c|c|c|c|}
\hline
\(P(K^2\ge k)\) & 0.050 & 0.010 & 0.001 \\
\hline
\(k\) & 3.841 & 6.635 & 10.828 \\
\hline
\end{tabular}
\end{center}

\end{problem}

\begin{solution}
（1）甲机床生产的产品中的一级品的频率为 \(\frac{150}{200}=75\%\).
乙机床生产的产品中的一级品的频率为 \(\frac{120}{200}=60\%\).

（2）由题中数据，\(K^2=\frac{400(150\times 80-120\times 50)^2}{270\times 130\times 200\times 200}=\frac{400}{39}>10>6.635\).
因此能有 \(99\%\) 的把握认为甲机床的产品质量与乙机床的产品质量有差异.
\end{solution}

\begin{problem}
记 \(S_n\) 为数列 \(\{a_n\}\) 的前 \(n\) 项和，已知 \(a_n>0\)，\(a_2=3a_1\)，且数列 \(\{\sqrt{S_n}\}\) 是等差数列，证明：\(\{a_n\}\) 是等差数列.

\end{problem}

\begin{solution}
因为数列 \(\{\sqrt{S_n}\}\) 是等差数列，设其公差为 \(d\)，则 \(d=\sqrt{S_2}-\sqrt{S_1}=\sqrt{a_1+a_2}-\sqrt{a_1}=\sqrt{4a_1}-\sqrt{a_1}=\sqrt{a_1}\).
所以 \(\sqrt{S_n}=\sqrt{a_1}+(n-1)\sqrt{a_1}=n\sqrt{a_1}\).
从而 \(S_n=a_1n^2\).
当 \(n\ge 2\) 时，\(a_n=S_n-S_{n-1}=a_1n^2-a_1(n-1)^2=2a_1n-a_1\).
当 \(n=1\) 时，\(2a_1\times 1-a_1=a_1\)，
也满足上式.
因此 \(a_n=2a_1n-a_1\).
于是 \(a_n-a_{n-1}=(2a_1n-a_1)-\bigl(2a_1(n-1)-a_1\bigr)=2a_1\).
故数列 \(\{a_n\}\) 是等差数列.
\end{solution}

\begin{problem}
已知直三棱柱 \(ABC-A_1B_1C_1\) 中，侧面 \(AA_1B_1B\) 为正方形，\(AB=BC=2\)，\(E\)，\(F\) 分别为 \(AC\) 和 \(CC_1\) 的中点，\(BF\perp A_1B_1\).

\begin{enumerate}
\item 求三棱锥 \(F-EBC\) 的体积；
\item 已知 \(D\) 为棱 \(A_1B_1\) 上的点，证明：\(BF\perp DE\).
\end{enumerate}

\end{problem}

\begin{solution}
（1）连结 \(AF\).
由题意可得 \(BF=\sqrt{BC^2+CF^2}=\sqrt{4+1}=\sqrt{5}\).
由 \(BF\perp A_1B_1\)，且 \(A_1B_1\parallel AB\)，得 \(AB\perp BF\).
又 \(AB\perp BB_1\)，且 \(BF\) 与 \(BB_1\) 是平面 \(BCC_1B_1\) 内相交的两条直线，所以 \(AB\perp \text{平面 }BCC_1B_1\)，从而 \(AB\perp BC\).
从而 \(AF=\sqrt{AB^2+BF^2}=\sqrt{4+5}=3\).
再由 \(AC=\sqrt{AF^2-CF^2}=\sqrt{9-1}=2\sqrt{2}\)，得 \(AB^2+BC^2=AC^2\)，所以 \(AB\perp BC\).
于是 \(\triangle ABC\) 为等腰直角三角形，\(S_{\triangle BCE}=\frac{1}{2}S_{\triangle ABC}=\frac{1}{2}\times \frac{1}{2}\times 2\times 2=1\).
因此 \(V_{F-EBC}=\frac{1}{3}S_{\triangle BCE}\times CF=\frac{1}{3}\times 1\times 1=\frac{1}{3}\).

（2）由（1）的结论，可将几何体补形为一个棱长为 \(2\) 的正方体 \(ABCM-A_1B_1C_1M_1\). 取棱 \(AM\)，\(BC\) 的中点分别为 \(H\)，\(G\)，连结 \(A_1H\)，\(HG\)，\(GB_1\).
在正方形 \(BCC_1B_1\) 中，\(G\)，\(F\) 为中点，所以 \(BF\perp B_1G\).
又已知 \(BF\perp A_1B_1\)，且 \(A_1B_1\cap B_1G=B_1\)，故 \(BF\perp \text{平面 }A_1B_1GH\).
因为 \(DE\subset \text{平面 }A_1B_1GH\)，所以 \(BF\perp DE\).
故证毕.
\end{solution}

\begin{problem}
设函数 \(f(x)=a^2x^2+ax-3\ln x+1\)，其中 \(a>0\).

\begin{enumerate}
\item 讨论 \(f(x)\) 的单调性；
\item 若 \(y=f(x)\) 的图象与 \(x\) 轴没有公共点，求 \(a\) 的取值范围.
\end{enumerate}

\end{problem}

\begin{solution}
（1）函数的定义域为 \((0,+\infty)\).
导函数为 \(f'(x)=2a^2x+a-\frac{3}{x}=\frac{2a^2x^2+ax-3}{x}=\frac{(2ax+3)(ax-1)}{x}\).
因为 \(a>0\)，\(x>0\)，所以 \(2ax+3>0\). 故
当 \(0<x<\frac{1}{a}\) 时，\(f'(x)<0\)；
当 \(x>\frac{1}{a}\) 时，\(f'(x)>0\).
所以 \(f(x)\) 的减区间为 \(\left(0,\frac{1}{a}\right)\)，增区间为 \(\left(\frac{1}{a},+\infty\right)\).

（2）因为
\(f(1)=a^2+a+1>0\)，
且 \(y=f(x)\) 的图象与 \(x\) 轴没有公共点，所以 \(y=f(x)\) 的图象在 \(x\) 轴上方.
由（1）可知，\(f(x)\) 在 \(x=\frac{1}{a}\) 处取得最小值，因此 \(f\left(\frac{1}{a}\right)>0\).
计算得 \(f\left(\frac{1}{a}\right)=a^2\cdot \frac{1}{a^2}+a\cdot \frac{1}{a}-3\ln \frac{1}{a}+1=1+1+3\ln a+1=3+3\ln a\).
于是 \(3+3\ln a>0\)，解得 \(a>\frac{1}{\e}\).
故 \(a\) 的取值范围为 \(\left(\frac{1}{\e},+\infty\right)\).
\end{solution}

\begin{problem}
抛物线 \(C\) 的顶点为坐标原点 \(O\)，焦点在 \(x\) 轴上，直线 \(l:x=1\) 交 \(C\) 于 \(P\)，\(Q\) 两点，且 \(OP\perp OQ\). 已知点 \(M(2,0)\)，且 \(\odot M\) 与 \(l\) 相切.

\begin{enumerate}
\item 求 \(C\)，\(\odot M\) 的方程；
\item 设 \(A_1\)，\(A_2\)，\(A_3\) 是 \(C\) 上的三个点，直线 \(A_1A_2\)，\(A_1A_3\) 均与 \(\odot M\) 相切. 判断直线 \(A_2A_3\) 与 \(\odot M\) 的位置关系，并说明理由.
\end{enumerate}

\end{problem}

\begin{solution}
（1）依题意设抛物线 \(C:y^2=2px\ (p>0)\)，又设 \(P(1,y_0)\)，\(Q(1,-y_0)\).
因为 \(OP\perp OQ\)，所以 \(\overrightarrow{OP}\cdot \overrightarrow{OQ}=1-y_0^2=0\).
又 \(P\)，\(Q\) 在抛物线上，故 \(y_0^2=2p\). 于是 \(1-2p=0\)，所以 \(2p=1\).
因此抛物线 \(C\) 的方程为 \(y^2=x\).
点 \(M(2,0)\) 到直线 \(x=1\) 的距离为 \(1\)，又 \(\odot M\) 与 \(x=1\) 相切，所以圆半径为 \(1\). 故 \(\odot M\) 的方程为 \((x-2)^2+y^2=1\).

（2）设 \(A_1(x_1,y_1)\)，\(A_2(x_2,y_2)\)，\(A_3(x_3,y_3)\).

先考虑 \(A_1A_2\) 斜率不存在的情形.

若 \(A_1A_2\) 的方程为 \(x=1\)，根据对称性，不妨设 \(A_1(1,1)\).
则过 \(A_1\) 与圆 \(M\) 相切的另一条直线方程为 \(y=1\). 此时该直线与抛物线只有一个交点，不存在 \(A_3\)，不合题意.

若 \(A_1A_2\) 的方程为 \(x=3\)，根据对称性，不妨设 \(A_1(3,\sqrt{3})\)，\(A_2(3,-\sqrt{3})\).
则过 \(A_1\) 与圆 \(M\) 相切的直线 \(A_1A_3\) 方程为 \(y-\sqrt{3}=\frac{\sqrt{3}}{3}(x-3)\).
又因为点 \(A_1\)，\(A_3\) 在抛物线 \(y^2=x\) 上，所以 \(k_{A_1A_3}=\frac{y_1-y_3}{x_1-x_3}=\frac{1}{y_1+y_3}=\frac{1}{\sqrt{3}+y_3}=\frac{\sqrt{3}}{3}\).
解得 \(y_3=0\)，从而 \(x_3=0\).
即 \(A_3(0,0)\).
此时直线 \(A_1A_3\) 与 \(A_2A_3\) 关于 \(x\) 轴对称，所以直线 \(A_2A_3\) 与圆 \(M\) 也相切.

若直线 \(A_2A_3\) 斜率不存在，则 \(x_2=x_3\)，从而 \(y_2=-y_3\). 由于直线 \(A_1A_2\)，\(A_1A_3\) 斜率存在，直线 \(A_1A_i\) 的方程为 \(x-(y_1+y_i)y+y_1y_i=0\)（\(i=2\)，\(3\)）.
由 \(A_1A_i\) 与圆 \(M\) 相切可得 \((y_1^2-1)y_i^2+2y_1y_i+3-y_1^2=0\)（\(i=2\)，\(3\)）.
若 \(y_1^2=1\)，上述方程至多有一个根，与两根互异矛盾，故 \(y_1^2\ne1\). 由韦达定理及 \(y_2+y_3=0\) 得 \(y_1=0\)，进而 \(y_2y_3=-3\). 所以 \(x_2=x_3=3\)，即直线 \(A_2A_3\) 为 \(x=3\)，与圆 \(M\) 相切.

下面考虑直线 \(A_1A_2\)，\(A_1A_3\)，\(A_2A_3\) 斜率都存在的情形.

因为 \(A_1\)，\(A_2\)，\(A_3\) 在抛物线 \(y^2=x\) 上，所以 \(x_1=y_1^2\)，\(x_2=y_2^2\)，\(x_3=y_3^2\).
于是 \(k_{A_1A_2}=\frac{y_1-y_2}{x_1-x_2}=\frac{1}{y_1+y_2}\)，\(k_{A_1A_3}=\frac{1}{y_1+y_3}\)，\(k_{A_2A_3}=\frac{1}{y_2+y_3}\).
因此直线 \(A_1A_2\) 的方程为
\(y-y_1=\frac{1}{y_1+y_2}(x-x_1)\)，整理得 \(x-(y_1+y_2)y+y_1y_2=0\).
由 \(y-y_1=\frac{1}{y_1+y_3}(x-x_1)\)，整理得 \(A_1A_3:\ x-(y_1+y_3)y+y_1y_3=0\).
由 \(y-y_2=\frac{1}{y_2+y_3}(x-x_2)\)，整理得 \(A_2A_3:\ x-(y_2+y_3)y+y_2y_3=0\).
因为直线 \(A_1A_2\) 与圆 \(M\) 相切，所以圆心 \(M(2,0)\) 到直线 \(A_1A_2\) 的距离为 \(1\)，即
\(\frac{|2+y_1y_2|}{\sqrt{1+(y_1+y_2)^2}}=1\).
整理得 \((y_1^2-1)y_2^2+2y_1y_2+3-y_1^2=0\).
因为直线 \(A_1A_3\) 与圆 \(M\) 相切，所以圆心 \(M(2,0)\) 到直线 \(A_1A_3\) 的距离为 \(1\)，即
\(\frac{|2+y_1y_3|}{\sqrt{1+(y_1+y_3)^2}}=1\).
整理得 \((y_1^2-1)y_3^2+2y_1y_3+3-y_1^2=0\).
所以 \(y_2\)，\(y_3\) 是方程
\((y_1^2-1)y^2+2y_1y+3-y_1^2=0\)
的两根，从而
\(y_2+y_3=-\frac{2y_1}{y_1^2-1}\)，\(y_2y_3=\frac{3-y_1^2}{y_1^2-1}\).
圆心 \(M\) 到直线 \(A_2A_3\) 的距离为
\examdisplaychain{
\frac{|2+y_2y_3|}{\sqrt{1+(y_2+y_3)^2}}
&=\frac{\left|2+\frac{3-y_1^2}{y_1^2-1}\right|}{\sqrt{1+\left(-\frac{2y_1}{y_1^2-1}\right)^2}}\\
&=\frac{\left|\frac{y_1^2+1}{y_1^2-1}\right|}{\sqrt{\frac{(y_1^2-1)^2+4y_1^2}{(y_1^2-1)^2}}}\\
&=\frac{|y_1^2+1|}{\sqrt{(y_1^2-1)^2+4y_1^2}}\\
&=1.
}
因此直线 \(A_2A_3\) 与圆 \(M\) 相切.

综上，直线 \(A_2A_3\) 与 \(\odot M\) 相切.
\end{solution}

\begin{problem}
在直角坐标系 \(xOy\) 中，以坐标原点为极点，\(x\) 轴正半轴为极轴建立极坐标系，曲线 \(C\) 的极坐标方程为 \(\rho=2\sqrt{2}\cos\theta\).

\begin{enumerate}
\item 将 \(C\) 的极坐标方程化为直角坐标方程；
\item 设点 \(A\) 的直角坐标为 \((1,0)\)，\(M\) 为 \(C\) 上的动点，点 \(P\) 满足 \(\overrightarrow{AP}=\sqrt{2}\overrightarrow{AM}\)，写出 \(P\) 的轨迹 \(C_1\) 的参数方程，并判断 \(C\) 与 \(C_1\) 是否有公共点.
\end{enumerate}

\end{problem}

\begin{solution}
（1）由曲线 \(C\) 的极坐标方程 \(\rho=2\sqrt{2}\cos\theta\) 可得 \(\rho^2=2\sqrt{2}\rho\cos\theta\).
将 \(x=\rho\cos\theta\)，\(y=\rho\sin\theta\) 代入，得 \(x^2+y^2=2\sqrt{2}x\)，即 \((x-\sqrt{2})^2+y^2=2\).

（2）设 \(P(x,y)\)，\(M(\sqrt{2}+\sqrt{2}\cos\theta,\sqrt{2}\sin\theta)\).
由 \(\overrightarrow{AP}=\sqrt{2}\overrightarrow{AM}\) 得 \((x-1,y)=\sqrt{2}\bigl(\sqrt{2}+\sqrt{2}\cos\theta-1,\sqrt{2}\sin\theta\bigr)\).
所以 \((x-1,y)=\bigl(2+2\cos\theta-\sqrt{2},2\sin\theta\bigr)\)，即 \(x=3-\sqrt{2}+2\cos\theta\)，\(y=2\sin\theta\)（\(\theta\) 为参数）.

因此 \(P\) 的轨迹 \(C_1\) 的参数方程为 \(x=3-\sqrt{2}+2\cos\theta\)，\(y=2\sin\theta\)（\(\theta\) 为参数）.
由（1）知，曲线 \(C\) 的圆心为 \((\sqrt{2},0)\)，半径为 \(\sqrt{2}\)；曲线 \(C_1\) 的圆心为 \((3-\sqrt{2},0)\)，半径为 \(2\).
两圆圆心距为 \((3-\sqrt{2})-\sqrt{2}=3-2\sqrt{2}\).
又 \(3-2\sqrt{2}<2-\sqrt{2}\)，
故两圆内含，没有公共点.
因此 \(C\) 与 \(C_1\) 没有公共点.
\end{solution}

\begin{problem}
已知函数 \(f(x)=|x-2|\)，\(g(x)=|2x+3|-|2x-1|\).

\begin{enumerate}
\item 画出 \(y=f(x)\) 和 \(y=g(x)\) 的图像；
\item 若 \(f(x+a)\ge g(x)\)，求 \(a\) 的取值范围.
\end{enumerate}

\end{problem}

\begin{solution}
（1）由绝对值的定义可得
\examdisplaycases{
f(x)=|x-2|=
}{
2-x, & x<2,\\
x-2, & x\ge 2
}
据此可作出 \(y=f(x)\) 的图像.

又
\examdisplaycases{
g(x)=|2x+3|-|2x-1|=
}{
-4, & x<-\frac{3}{2},\\
4x+2, & -\frac{3}{2}\le x<\frac{1}{2},\\
4, & x\ge \frac{1}{2}
}
据此可作出 \(y=f(x)\)，\(y=g(x)\) 以及第（2）问中临界位置时 \(y=f(x+a)\) 的示意图如下：

\bitmapfigure[max width=5.8cm,max totalheight=3.8cm]{img/2021/national_paper_A_liberal/q23_solution_fig1.png}

（2）由 \(f(x+a)=|x+a-2|\)
可知，\(y=f(x+a)\) 是将 \(y=f(x)\) 平移了 \(a\) 个单位得到的图像.
由 \(x=\frac{1}{2}\) 处的条件，得
\(\left|\frac{1}{2}+a-2\right|\ge 4\)，即 \(a\le-\frac{5}{2}\) 或 \(a\ge\frac{11}{2}\).
若 \(a\le-\frac{5}{2}\)，则 \(x_0=2-a\ge\frac{1}{2}\)，从而
\(f(x_0+a)=0<g(x_0)=4\)，不合题意.
因此只需考虑 \(a\ge\frac{11}{2}\). 此时当 \(-\frac{3}{2}\le x<\frac{1}{2}\) 时，
\(f(x+a)-g(x)=x+a-2-(4x+2)=a-4-3x\ge a-\frac{11}{2}\ge0\)；
当 \(x\ge\frac{1}{2}\) 时，\(f(x+a)=x+a-2\ge4=g(x)\)；
当 \(x< -\frac{3}{2}\) 时，\(g(x)=-4\)，不等式显然成立.
所以 \(a\) 的取值范围为
 \(\left[\frac{11}{2},+\infty\right)\).
\end{solution}
